\newpage
\phantomsection
\label{prf:corestriction-surjective}
\LRAProofFor{prop:corestriction-surjective}

\begin{remark*}[Return]
\hyperref[prop:corestriction-surjective]{Return to Proposition}
\end{remark*}

\begin{proposition*}[Corestriction Is Surjective]
For every function $f:A\to B$, the corestriction
$f_{\operatorname{im}}:A\to\im(f)$ is surjective.
\end{proposition*}

\begin{proof}[Professional Standard Proof]
\LRAProofBodyStart
Let $y\in\im(f)$. By definition of image, there is $a\in A$ such that
$f(a)=y$. Since $f_{\operatorname{im}}(a)=f(a)$, we have
$f_{\operatorname{im}}(a)=y$. Hence every element of $\im(f)$ is hit.
\end{proof}

\begin{proof}[Detailed Learning Proof]
\LRAProofBodyStart
To prove surjectivity, start with an arbitrary element of the codomain
$\im(f)$. Such an element is in the image exactly when it is equal to $f(a)$
for some input $a$. The corestricted function has the same value as $f$ at
that input, so it reaches the chosen element.
\end{proof}
\begin{remark*}[Proof structure]
\end{remark*}



\begin{dependencies}
\begin{itemize}
  \item \hyperref[prop:corestriction-surjective]{Proof target}
\end{itemize}
\end{dependencies}

\clearpage

